\documentclass[11pt,a4paper]{report}
\usepackage[utf8]{inputenc}
\usepackage[T1]{fontenc}
\usepackage{lmodern,microtype}
\usepackage[margin=2.4cm]{geometry}
\usepackage{graphicx}
\usepackage{booktabs,longtable,array,amsmath,amssymb,xcolor,listings,hyperref,enumitem,titlesec,etoolbox}
\hypersetup{colorlinks=true,linkcolor=blue!50!black,urlcolor=blue!50!black,citecolor=blue!50!black,pdftitle={nwalaam: Autologistic actor attribute models in Stata. Manual},pdfauthor={Thomas U. Grund}}
\lstdefinestyle{stata}{basicstyle=\ttfamily\footnotesize,backgroundcolor=\color{black!4},frame=single,rulecolor=\color{black!20},breaklines=true,columns=fullflexible,keepspaces=true,xleftmargin=2pt,xrightmargin=2pt,aboveskip=8pt,belowskip=8pt}
\lstset{style=stata}
\setlength{\emergencystretch}{3em}
\newcommand{\nwalaam}{\texttt{nwalaam}}
\newcommand{\alaamee}{\texttt{ALAAMEE}}
\newcommand{\mpnet}{\texttt{MPNet}}
% \opt sets an option or effect name in typewriter type and lets long names break after ( : , _ and / so that they do not run into the next table column
\ExplSyntaxOn
\tl_new:N \l__nwalaam_opt_tl
\NewDocumentCommand{\opt}{m}
  {
    \tl_set:Nn \l__nwalaam_opt_tl {#1}
    \regex_replace_all:nnN { ([\:\,]) } { \1 \c{allowbreak}\cB\{\cE\} } \l__nwalaam_opt_tl
    \regex_replace_all:nnN { \c{\_} } { \c{\_} \c{allowbreak}\cB\{\cE\} } \l__nwalaam_opt_tl
    \regex_replace_all:nnN { ([a-z])([A-Z]) } { \1 \c{allowbreak}\cB\{\cE\} \2 } \l__nwalaam_opt_tl
    \regex_replace_all:nnN { (/) } { \1 \c{allowbreak}\cB\{\cE\} } \l__nwalaam_opt_tl
    \regex_replace_all:nnN { ([\|\(\[=]) } { \1 \c{allowbreak}\cB\{\cE\} } \l__nwalaam_opt_tl
    \texttt{ \tl_use:N \l__nwalaam_opt_tl }
  }
\ExplSyntaxOff
% the listing of the output of an example (manual/ex/out/<name>.txt, written by the do-files of manual/ex)
\newcommand{\showlog}[1]{\lstinputlisting{ex/out/#1.txt}}
\newcommand{\inc}[1]{\typeout{>>>#1}\input{#1}}
\newcommand{\manualseries}{1.4}
% The manual is labelled by the minor series (valid for 1.4.x): change \manualseries only for a new minor or major release whose documented behaviour changed.
\newcommand{\pkgversion}{1.4.0}
\newcommand{\pkgdate}{22 November 2026}
% Exact version and date of the release that produced the listings (written by tools/set_version.py; rerun the examples when it changes).
% DOI of the archived manual (Zenodo). PLACEHOLDER: no record exists yet; put the all-versions DOI of the new record here (for example 10.5281/zenodo.NNNNNNN) after the upload and rebuild. Nothing is printed for an empty DOI.
\newcommand{\manualdoi}{10.5281/zenodo.24397901}% DOI of the Zenodo record
\newcommand{\citedoi}{\ifdefempty{\manualdoi}{}{ Zenodo. \url{https://doi.org/\manualdoi}}}
\title{\textbf{nwalaam}\\[4pt]\large Autologistic actor attribute models in Stata\\ (a package that builds on \texttt{nwcommands})\\[6pt]\normalsize User manual and reference}
\author{Thomas U. Grund\\ RWTH Aachen University}
\date{This manual describes \texttt{nwalaam} \manualseries.x\\ (examples produced with \texttt{nwalaam} \pkgversion, \pkgdate)}
\begin{document}
\maketitle
\begin{quote}\noindent\textbf{How to cite.} This manual describes the series~\manualseries{} of \nwalaam{} (it holds for all versions \manualseries.x); the listings were produced with version~\pkgversion{} (\pkgdate). It is the citable reference for the package. Please cite it as:\par\smallskip
\noindent Grund, T.~U. (2026). \emph{nwalaam: Autologistic actor attribute models in Stata. Manual.}\citedoi\par\smallskip
\noindent Please also report the exact version of the program that produced your results, as printed by \texttt{nwalaam version} and stored in \texttt{e(version)}.
\end{quote}
\chapter*{Acknowledgments}\addcontentsline{toc}{chapter}{Acknowledgments}
\nwalaam{} follows the software \alaamee{} of Alex Stivala, Peng Wang and Alessandro Lomi~\cite{stivala2024} and exists because of it. The autologistic actor attribute model is due to Robins, Pattison and Elliott~\cite{robins2001} and Daraganova and Robins~\cite{daraganova2013}; \mpnet{}~\cite{wang2009} is the other reference implementation. The equilibrium-expectation algorithm is that of Borisenko, Byshkin and Lomi~\cite{borisenko2019}, the stochastic approximation that of Snijders~\cite{snijders2002}. The definitions of the effects, the conditional sampler for snowball data and the zero-or-one sampler for rare outcomes are those documented by Stivala and co-authors~\cite{stivala2020,stivala2023,stivala2024}. The Stata code was written independently from these publications; no code of \alaamee{} is included, and its change statistics are used as the reference against which the code was checked.

\nwalaam{} is not affiliated with or endorsed by the authors of \alaamee{} or \mpnet{}, and its errors are the author's. It is meant for researchers who work in Stata or prefer it to R or Python. It is not a replacement for \alaamee{} and \mpnet{}, which remain the reference implementations.

\paragraph{Use of an AI coding assistant.} The software was developed with Claude Code (Anthropic), an AI coding assistant, working under the author's direction. It helped port the algorithms into Stata, Mata and C and helped write the tests. The author specified the scope, reviewed the code, and is responsible for the package. Its correctness rests on validation, not on the assistant: the change statistics are compared numerically with those of \alaamee{}, the estimators with the exact maximum likelihood estimate where it can be computed, and the test suite reruns these comparisons.

\tableofcontents
\inc{ch01_intro}
\inc{ch02_model}
\inc{ch03_start}
\inc{ch04_options}
\inc{ch05_effects}
\inc{ch06_estimation}
\inc{ch07_special}
\inc{ch08_postest}
\inc{ch09_examples}
\inc{ch10_validation}
\inc{ch11_limits}
\appendix
\inc{app_translation}
\bibliographystyle{plain}
\begin{thebibliography}{99}
\bibitem{robins2001} Robins, G., Pattison, P., Elliott, P. (2001). Network models for social influence processes. \emph{Psychometrika} 66(2), 161--189. \url{https://doi.org/10.1007/BF02294834}
\bibitem{daraganova2013} Daraganova, G., Robins, G. (2013). Autologistic actor attribute models. In: Lusher, D., Koskinen, J., Robins, G. (eds.), \emph{Exponential Random Graph Models for Social Networks}, chapter 9, 102--114. Cambridge: Cambridge University Press.
\bibitem{stivala2024} Stivala, A., Wang, P., Lomi, A. (2024). ALAAMEE: Open-source software for fitting autologistic actor attribute models. \emph{PLOS Complex Systems} 1(4), e0000021.
\bibitem{borisenko2019} Borisenko, A., Byshkin, M., Lomi, A. (2019). A simple algorithm for scalable Monte Carlo inference. arXiv:1901.00533.
\bibitem{snijders2002} Snijders, T.A.B. (2002). Markov chain Monte Carlo estimation of exponential random graph models. \emph{Journal of Social Structure} 3(2), 1--40.
\bibitem{stivala2020} Stivala, A., Gallagher, H.C., Rolls, D.A., Wang, P., Robins, G.L. (2020). Using sampled network data with the autologistic actor attribute model. arXiv:2002.00849.
\bibitem{stivala2023} Stivala, A. (2023). Overcoming near-degeneracy in the autologistic actor attribute model. arXiv:2309.07338.
\bibitem{wang2009} Wang, P., Robins, G., Pattison, P. (2009). \emph{PNet: A program for the simulation and estimation of exponential random graph (p*) models.} Melbourne: Melbourne School of Psychological Sciences, The University of Melbourne.
\bibitem{zachary1977} Zachary, W.W. (1977). An information flow model for conflict and fission in small groups. \emph{Journal of Anthropological Research} 33(4), 452--473.
\bibitem{lazega2001} Lazega, E. (2001). \emph{The Collegial Phenomenon: The Social Mechanisms of Cooperation Among Peers in a Corporate Law Partnership}. Oxford: Oxford University Press.
\bibitem{michell1996} Michell, L., West, P. (1996). Peer pressure to smoke: the meaning depends on the method. \emph{Health Education Research} 11(1), 39--49.
\bibitem{davis1941} Davis, A., Gardner, B.B., Gardner, M.R. (1941). \emph{Deep South: A Social Anthropological Study of Caste and Class}. Chicago: University of Chicago Press.
\bibitem{manski1993} Manski, C.F. (1993). Identification of endogenous social effects: the reflection problem. \emph{Review of Economic Studies} 60(3), 531--542.
\end{thebibliography}
\end{document}
